\documentclass[10pt,a4paper]{amsart}
\usepackage[margin=19mm]{geometry}
\usepackage{amsmath,amssymb,mathtools}
\usepackage[scr=rsfs,cal=euler]{mathalfa}
\usepackage{xfrac}
\usepackage[unicode,hidelinks,bookmarks=false]{hyperref}
\newcommand{\ie}{i.e.\ }
\newcommand{\cf}{cf.\ }
\newcommand{\resp}{respectively\ }
\newcommand{\Subsection}{\subsection}
\providecommand{\nobreakdash}{\nobreak\hbox{-}}
\setlength{\emergencystretch}{2em}
\multlinegap0pt
\usepackage{amssymb}
\newtheorem{theorem}{Theorem}
\newtheorem{lemma}{Lemma}
\newtheorem{proposition}{Proposition}
\title{Shadowing Endomorphisms of Compact Groups}
\author{Dekui Peng}
\date{Source: arXiv:2608.00955v1 (CC BY 4.0)}
\begin{document}
\maketitle

\noindent\emph{Source and licence.} Shadowing Endomorphisms of Compact Groups. Version arXiv:2608.00955v1 (CC BY 4.0), distributed by arXiv under the Creative Commons Attribution 4.0 International licence, http://creativecommons.org/licenses/by/4.0/, as recorded in the arXiv licence metadata for this exact version and shown on the abstract page https://arxiv.org/abs/2608.00955v1. Full licence notice: \texttt{licenses/arxiv-2608.00955-CC-BY-4.0.txt}.\par\smallskip

\noindent\textit{Editorial scope.} Counted unit: the article's theorem thm:abelian-endo-shad (source0.tex lines 1618--1627). The article's complete original proof of it is delivered with it (source0.tex lines 1629--1827, one \texttt{proof} environment of 199 lines). No mathematics is written, repaired or re-derived: the statement and the proof are the article's own lines, in the article's own notation, and no retained line is substituted or re-flowed.  Retained uncounted as the source-local context the proof needs: lines 537--540; lines 1016--1021; lines 1059--1066; lines 1493--1497.  Nothing was omitted from any retained span, so the package carries no omission record.  No retained line contains a citation, so the wrapper carries no bibliography and no bibliography entry.  No retained line contains a figure environment, an image-inclusion command or drawing code, so no image is delivered and the package declares no asset dependency.  Editorial formatting only, in addition: an AMS \texttt{amsart} preamble that restates the article's own declarations and macros used by the retained lines, namely the environments \texttt{theorem}, \texttt{lemma}, \texttt{proposition} on separate counters, together with the standard packages those lines need.  The retained environments print in this document as Theorem 1, Lemma 1, Proposition 1 in order of appearance, whereas the article prints the same items with the numbering of its own full document.  The complete native source of the article as distributed in the arXiv e-print for this exact version is bundled unchanged as \texttt{sources/206581/source0.tex}.
\begin{theorem}[Inverse-limit Shadowing Theorem]\label{thm:inverse-limit-shadowing}
 Let $(I,\leq)$ be a directed set, and let \[ \mathbf X = \bigl\{(X_i,f_i),p_i^j:i\leq j\text{ in }I\bigr\} \] be an inverse system of compact dynamical systems satisfying the Mittag--Leffler condition. 
 Suppose that $f_i$ has shadowing for every $i\in I$. Then the induced map \[ f:\varprojlim_{i\in I}X_i\longrightarrow \varprojlim_{i\in I}X_i \] has shadowing. 
 \end{theorem}

\begin{lemma}\label{lem:quotient-shadowing}
Let $G$ and $H$ be topological groups, let $\alpha:G\to G$ and
$\beta:H\to H$ be endomorphisms, and let $\pi:G\to H$ be an open continuous surjective homomorphism satisfying
$\pi\circ\alpha=\beta\circ\pi$. If $\alpha$ has finite shadowing, then
$\beta$ has finite shadowing.
\end{lemma}

\begin{lemma}\label{lem:extension-shadowing}
Let $G$ be a compact group, let $\alpha:G\to G$ be an endomorphism,
and let $N\trianglelefteq G$ be a closed normal subgroup such that
$\alpha(N)\subseteq N$. Let $\alpha_N:=\alpha|_N$, and let
$\bar\alpha:G/N\to G/N$ be the induced endomorphism. If both
$\alpha_N$ and $\bar\alpha$ have finite shadowing, then $\alpha$ has
finite shadowing.
\end{lemma}

\begin{proposition}\label{prop:finite-dim}
If $A$ is finite-dimensional, or equivalently, if $\Gamma$ is of
finite rank, then $\beta$ has shadowing if and only if
$T_{\mathbb Q}$ is hyperbolic.
\end{proposition}

\begin{theorem}\label{thm:abelian-endo-shad}
Let $A$ be a connected compact abelian group and let
$\beta:A\to A$ be an endomorphism. Then the following conditions are
equivalent:
\begin{enumerate}
\item $\beta$ has shadowing;
\item $T_{\mathbb Q}|_W$ is hyperbolic for every finite-dimensional
$T_{\mathbb Q}$-invariant subspace $W\leq V$.
\end{enumerate}
\end{theorem}

\begin{proof}
Assume first that $\beta$ has shadowing, and let $W\leq V$ be a
finite-dimensional $T_{\mathbb Q}$-invariant subspace. Put
$\Delta:=W\cap\Gamma$. Then $\Delta$ is a $T$-invariant subgroup of
$\Gamma$. Since $W$ is a finite-dimensional rational subspace of
$\Gamma\otimes_{\mathbb Z}\mathbb Q$, the subgroup
$\Delta=W\cap\Gamma$ has finite rank and spans $W$ over $\mathbb Q$.
Thus
$
\Delta\otimes_{\mathbb Z}\mathbb Q=W.
$

The inclusion $\Delta\hookrightarrow\Gamma$ dualizes to an equivariant
surjective homomorphism $A\to\widehat\Delta$. Hence the endomorphism
induced by $\beta$ on the finite-dimensional connected compact abelian group
$\widehat\Delta$ has shadowing. By Proposition \ref{prop:finite-dim}, $T_{\mathbb Q}|_W$ has no eigenvalue on the unit circle.

Conversely, assume that every finite-dimensional
$T_{\mathbb Q}$-invariant subspace of $V$ is hyperbolic. Regard $V$ as
a module over $\mathbb Q[u]$ by putting $u\cdot v=T_{\mathbb Q}v$.
Define its algebraic part by
\begin{equation*}
V_{\mathrm{alg}}
:=
\{v\in V:p(T_{\mathbb Q})v=0
\text{ for some }0\neq p\in\mathbb Q[u]\},
\end{equation*}
and put $\Gamma_{\mathrm{alg}}:=\Gamma\cap V_{\mathrm{alg}}$.

We have
\[
\Gamma_{\mathrm{alg}}\otimes_{\mathbb Z}\mathbb Q
=
V_{\mathrm{alg}}.
\]
Indeed, let $v\in V_{\mathrm{alg}}$. Choose a nonzero integer $m$
such that $mv\in\Gamma$. If
$p(T_{\mathbb Q})v=0$ for some nonzero $p\in\mathbb Q[u]$, then,
after clearing denominators,
$q(T)(mv)=0$ for some nonzero $q\in\mathbb Z[u]$. Hence
$mv\in\Gamma_{\mathrm{alg}}$, and therefore
$v\in\Gamma_{\mathrm{alg}}\otimes_{\mathbb Z}\mathbb Q$. The other inclusion is obvious.

We now consider the system dual to $\Gamma_{\mathrm{alg}}$. Let $\mathcal F$ be the family of all finite subsets of $\Gamma_{\mathrm{alg}}$ and for each $F\in \mathcal F$, let $W_F$ be the $\mathbb Q[u]$-submodule of
$V_{\mathrm{alg}}$ generated by $F$, and
put $\Delta_F:=W_F\cap\Gamma_{\mathrm{alg}}$.

Since every generator of $W_F$ is annihilated by a nonzero polynomial,
$W_F$ is finite-dimensional over $\mathbb Q$. The induced
endomorphism on $\widehat{\Delta_F}$ therefore has shadowing by the
finite-dimensional solenoidal criterion Proposition \ref{prop:finite-dim} and the hypothesis on the
spectrum.

Order $\mathcal F$ by inclusion. If $F\subseteq F'$, then
$\Delta_F\subseteq\Delta_{F'}$, and the inclusion dualizes to a
surjective equivariant homomorphism
$
\widehat{\Delta_{F'}}
\longrightarrow
\widehat{\Delta_F}.
$
Since
$
\Gamma_{\mathrm{alg}}
=
\bigcup_{F\in\mathcal F}\Delta_F,
$
Pontryagin duality gives
$
\widehat{\Gamma_{\mathrm{alg}}}
\cong
\varprojlim_{F\in\mathcal F}\widehat{\Delta_F}.
$
All bonding maps are surjective. Hence
Theorem~\ref{thm:inverse-limit-shadowing} implies that the induced
endomorphism on $\widehat{\Gamma_{\mathrm{alg}}}$ has shadowing.

Now put $\Lambda:=\Gamma/\Gamma_{\mathrm{alg}}$. Its rationalization
is naturally isomorphic to $V/V_{\mathrm{alg}}$, which is torsion-free
as a $\mathbb Q[u]$-module. Indeed, if
$p(T_{\mathbb Q})(v+V_{\mathrm{alg}})=0$, then
$p(T_{\mathbb Q})v\in V_{\mathrm{alg}}$, and hence
$q(T_{\mathbb Q})\big(p(T_{\mathbb Q})v\big)=0$ for some nonzero
$q\in\mathbb Q[u]$. Thus $v\in V_{\mathrm{alg}}$.

Again, we let $\mathcal E$ be the set of finite subsets of $\Lambda$ and for each $E\in \mathcal E$, let $M_E$ be the $\mathbb Q[u]$-submodule of
$\Lambda\otimes\mathbb Q$ generated by $E$, and put
$\Lambda_E:=M_E\cap\Lambda$. Since $\mathbb Q[u]$ is a principal ideal
domain and $M_E$ is finitely generated and torsion-free, there is an
integer $r_E$ such that
\[
M_E\cong\mathbb Q[u]^{r_E}.
\]

The dual of multiplication by $u$ on $\mathbb Q[u]^{r_E}$ is a
one-sided full shift on
$(\widehat{\mathbb Q}^{\,r_E})^{\mathbb N}$, and hence has shadowing.
Indeed, as an abelian group, one has
\[
\mathbb Q[u]^{r_E}
=
\bigoplus_{j\geq 0}\mathbb Q^{r_E}u^j.
\]
Therefore, by Pontryagin duality,
\[
\widehat{\mathbb Q[u]^{r_E}}
\cong
\prod_{j\geq 0}\widehat{\mathbb Q^{r_E}}
\cong
\bigl(\widehat{\mathbb Q}^{\,r_E}\bigr)^{\mathbb N}.
\]
More explicitly, every element of $\mathbb Q[u]^{r_E}$ can be written
uniquely in the form
\[
p(u)=\sum_{j\geq 0}a_j u^j,
\qquad a_j\in\mathbb Q^{r_E},
\]
where all but finitely many $a_j$ are zero. Under the above
identification, a point
$x=(x_0,x_1,\ldots)\in
(\widehat{\mathbb Q}^{\,r_E})^{\mathbb N}$
corresponds to the character $\chi_x$ defined by
\[
\chi_x(p)
=
\prod_{j\geq 0}x_j(a_j)\in \mathbb T.
\]

Let $m_u:\mathbb Q[u]^{r_E}\to\mathbb Q[u]^{r_E}$ denote
multiplication by $u$. Its dual map satisfies
\[
\widehat{m_u}(\chi_x)(p)
=
\chi_x(up).
\]
Since
$
up(u)=\sum_{j\geq 0}a_j u^{j+1},
$
we obtain
$
\chi_x(up)
=
\prod_{j\geq 0}x_{j+1}(a_j).
$
Thus, under the identification above,
\[
\widehat{m_u}(x_0,x_1,x_2,\ldots)
=
(x_1,x_2,x_3,\ldots).
\]
Hence the dual of multiplication by $u$ on
$\mathbb Q[u]^{r_E}$ is precisely the one-sided full shift on
$\bigl(\widehat{\mathbb Q}^{\,r_E}\bigr)^{\mathbb N}$.



The inclusion $\Lambda_E\hookrightarrow M_E$ dualizes to an
equivariant surjective homomorphism
$\widehat{M_E}\to\widehat{\Lambda_E}$. It follows from Lemma \ref{lem:quotient-shadowing} that the induced
endomorphism on $\widehat{\Lambda_E}$ has shadowing.



Order $\mathcal E$ by inclusion. If $E\subseteq E'$, then
$M_E\subseteq M_{E'}$ and hence
$\Lambda_E\subseteq\Lambda_{E'}$. Moreover,
$
\Lambda=\bigcup_{E\in\mathcal E}\Lambda_E.
$
Therefore
$
\widehat\Lambda
\cong
\varprojlim_{E\in\mathcal E}\widehat{\Lambda_E},
$
where all bonding maps are surjective and equivariant. Theorem~
\ref{thm:inverse-limit-shadowing} now implies that the induced
endomorphism on $\widehat\Lambda$ has shadowing.

Finally, the exact sequence
\[
0\longrightarrow\Gamma_{\mathrm{alg}}
\longrightarrow\Gamma
\longrightarrow\Lambda
\longrightarrow 0
\]
dualizes to an equivariant exact sequence
\[
1\longrightarrow\widehat\Lambda
\longrightarrow A
\longrightarrow\widehat{\Gamma_{\mathrm{alg}}}
\longrightarrow 1.
\]
Both the restriction to $\widehat\Lambda$ and the induced
endomorphism on $\widehat{\Gamma_{\mathrm{alg}}}$ have shadowing.
Lemma~\ref{lem:extension-shadowing} therefore implies that $\beta$ has
shadowing.
\end{proof}

\end{document}
